\section{Introducción}

La planificación por lotería (\textit{lottery scheduling}) representa
derechos de uso del CPU mediante boletos: en cada decisión de despacho se
sortea un boleto al azar entre las tareas activas y la propietaria del
boleto ganador recibe el siguiente intervalo de ejecución. A largo plazo, si
la selección es uniforme, la fracción de CPU que recibe una tarea converge a
su fracción de los boletos activos --- la propiedad de \textit{proportional
share} que motiva el mecanismo, descrita originalmente por Waldspurger y
Weihl~\cite{Waldspurger1994}.

Este proyecto traduce esa política a un mecanismo ejecutable y verificable:
un planificador en espacio de usuario, escrito en C17 sobre \textit{pthreads}
POSIX, que no modifica el \textit{kernel} ni sustituye el planificador de
GNU/Linux. El programa simula la existencia de un único CPU lógico ---cada
tarea es un hilo que solo avanza su trabajo cuando el planificador se lo
permite explícitamente--- y usa como carga de trabajo instrumentada el
cálculo incremental de la serie de Taylor de $\arcsin(1) = \pi/2$, elegida
por tener estado real que preservar entre activaciones (término anterior,
suma acumulada, índice) y por ser verificable contra un valor conocido.

\subsection{Alcance de este informe}

Las secciones siguientes cubren, en orden: la arquitectura del sistema y el
protocolo de sincronización (datos compartidos, mutex, variables de
condición y sus predicados de espera); el pseudocódigo de los algoritmos
centrales del sorteo, la activación, la cesión y la terminación; la
verificación contra los siete casos mínimos de funcionamiento que exige el
enunciado, incluyendo la evidencia de \textit{sanitizers}; los resultados del
experimento de proporcionalidad; el diseño e implementación de la extensión
elegida (\textit{compensation tickets}); un análisis comparativo con el paper
original de Waldspurger y Weihl, incluyendo las diferencias con su prototipo
sobre Mach 3.0 y las limitaciones de este prototipo en espacio de usuario; y
una discusión de los \textit{trade-offs} de diseño más relevantes. El código
fuente completo, la trazabilidad del equipo y el historial de decisiones de
diseño extendido viven en los repositorios públicos del proyecto, citados
donde corresponde.
